\section{Transformer 的跨域应用：同一计算模式，不同数据接口}

课堂用七页展示 Transformer 从文本扩展到语音、音乐、视觉、机器人、游戏和生物。真正可迁移的结论不是“所有任务都被统一”，而是许多领域都能把对象表示成 token 或 set，再用 attention 建模远距离交互。数据表示、输出约束和评估标准仍然高度领域化。

\subsection{文本与语言：最成熟的 token interface}

文本天然是离散序列，因而最容易接入 autoregressive objective。翻译、摘要、问答、代码和对话看似不同，往往都可转写为条件 token generation。下面的应用页应按输入、输出与验证标准拆解，不能因为都输出文字就认为任务结构相同。

\lecturefigure{slide-045.jpg}{Transformer 在文本生成、理解与对话中的应用}{官方 CS25 V4 slide deck，第 45 页}

\begin{knowledgebox}{读图：统一接口不等于统一任务}
把任务写成 text-to-text 可以共享模型与训练代码，但 loss、decoding、事实性和评价目标仍不同。翻译关注语义保持，摘要关注覆盖与压缩，代码关注可执行性，对话还要处理上下文、人格与安全规范。
\end{knowledgebox}

\subsection{Audio：波形需要先变成可建模表示}

音频是高频连续信号。系统通常先把波形转换为 spectrogram、codec token 或 learned feature，再用 Transformer 建模时间关系；输出也可能是文字、声学 token 或音乐事件。读图时先找 signal representation，再判断模型是在识别、生成还是跨模态对齐。

\lecturefigure{slide-046.jpg}{Speech 与 Music Transformer 的应用图景}{官方 CS25 V4 slide deck，第 46 页}

\teachervoice{这组 application slide 只做方向导览，讲者把深入机制留给后续嘉宾。讲义因此解释共同接口和领域差异，不把每个模型名称扩写成未经课堂支持的产品结论。}

\begin{importantbox}{模态建模的共同三步}
\begin{enumerate}
  \item \textbf{Tokenizer/encoder}：把原始信号变成有限长度表示；
  \item \textbf{Backbone}：用 attention 建模表示之间的条件关系；
  \item \textbf{Decoder/rendering}：把 hidden state 转回文字、图像、声音或动作。
\end{enumerate}
统一的是中间计算，最难的往往是两端的表示与评价。
\end{importantbox}

\subsection{Vision：patch、frame 与跨模态对齐}

Vision Transformer（视觉 Transformer，ViT）把图像切成 patch token；视频还要编码时间维。分析任务输出 label、box 或文字，生成任务则需要把 latent/token 解码成像素序列。接下来的两页分别对应 perception 与 generation，比较重点是 tokenization、decoder 和评价目标。

\lecturefigure{slide-047.jpg}{Vision Transformer 用 patch token 分析图像与视频}{官方 CS25 V4 slide deck，第 47 页}

\begin{importantbox}{Patch embedding}
若图像大小为 $H\times W$，patch 大小为 $P\times P$，则 token 数约为
\[
N=\frac{HW}{P^2}.
\]
更小 patch 保留更多局部细节，却让 attention 的 $N^2$ 成本迅速增加。视觉系统因此常使用分层结构、局部 attention、downsampling 或 latent bottleneck。
\end{importantbox}

\lecturefigure{slide-048.jpg}{Diffusion 与 Transformer 支持图像、视频生成}{官方 CS25 V4 slide deck，第 48 页}

\begin{warningbox}{生成画面真实不等于世界模型正确}
图像和视频模型可以产生高保真纹理，却仍可能违反物体恒常、几何、因果或长时间一致性。视觉质量 metric 与物理正确性不是同一目标，不能用“看起来像”替代对结构和动态的验证。
\end{warningbox}

\subsection{Robotics、游戏与生物：输出必须进入外部约束}

机器人和游戏把输出映射到 action；生物模型把序列或三维结构映射到功能与预测。此时错误不再只是生成一个坏句子，而可能改变物理轨迹、策略或实验选择。下面三页要从环境反馈和验证成本出发阅读，而不是把跨域模型只当成应用拼图。

\lecturefigure{slide-049.jpg}{机器人、仿真与 physical task 中的 Transformer}{官方 CS25 V4 slide deck，第 49 页}

\begin{knowledgebox}{术语消化：policy 与 world feedback}
Policy（策略）把 observation 映射到 action。语言模型常一次生成 token，机器人 policy 则持续接收传感器反馈并更新动作。闭环中的延迟、误差累积、碰撞约束和恢复机制，使 embodied system 比离线生成更难验收。
\end{knowledgebox}

\lecturefigure{slide-050.jpg}{Transformer 在围棋、实时策略与游戏环境中的应用}{官方 CS25 V4 slide deck，第 50 页}

\lecturefigure{slide-051.jpg}{医疗语言模型、蛋白结构与生物序列应用}{官方 CS25 V4 slide deck，第 51 页}

\begin{warningbox}{领域名称列表必须回到任务与证据}
Med-PaLM、AlphaFold、robotics transformer 与 game agent 解决的对象、监督信号和安全要求完全不同。记住名称不如记住三件事：输入如何表示、输出如何验证、错误由谁承担。
\end{warningbox}

\subsection{本章小结}

Transformer 的跨域扩张来自 attention 对集合与序列关系的通用建模能力，但每个领域仍需要自己的 tokenizer、decoder、constraint 与 evaluation。统一 backbone 降低了模型设计成本，没有消除领域知识和闭环验证。

